\section{Exact Gallery-Level Rank-Retention Geometry}
\label{app:ranktheory}
This section gives an exact distance to failure for a fixed retrieval problem.
It also gives a complementary score bound and an aggregate KL budget bound.
None changes the empirical selection rule or primary R@1 endpoints.
Surrogate calibration already connects surrogate losses with decision errors \citep{bartlett2003risk,pires2016calibration}.
Here the specific result is the exact positive-set failure radius for a fixed retrieval gallery.
The proof uses standard convex projection arguments.

\subsection{Definitions}
Fix one query and a finite gallery with $n$ candidates.
Let $R$ be its nonempty relevant set and $N$ its irrelevant set.
Teacher and student distributions $q,p$ lie in the same probability simplex.
Define forward KL in nats by
\begin{equation}
D(q\|p)=\sum_{i=1}^n q_i\log\frac{q_i}{p_i}.
\end{equation}
Use $0\log(0/p)=0$.
A positive numerator with zero denominator contributes infinity.
For softmax probabilities, both models use the same gallery and positive temperature.

For an integer $K\geq1$, define the closed failure event
\begin{equation}
F_K=\left\{p:\left|\{j\in N:p_j\geq\max_{i\in R}p_i\}\right|\geq K\right\}.
\end{equation}
Cross-relevance ties count against the relevant candidate.
Outside $F_K$, every tie rule places a relevant item within the first $K$ positions.
If $|N|<K$, failure is impossible.

Write $a_1\geq\cdots\geq a_r$ for relevant teacher probabilities.
Write $b_1\geq\cdots\geq b_s$ for irrelevant teacher probabilities.
When $K\leq s$, strict teacher correctness means $a_1>b_K$.
The exact distance to failure is
\begin{equation}
C_K(q,R)=\min_{p\in F_K}D(q\|p).
\end{equation}

\subsection{Exact positive-set projection}
\begin{theorem}
\label{thm:exactprojection}
If $|N|<K$, then $C_K=+\infty$.
If $a_1\leq b_K$, then $C_K=0$.
Otherwise, let $S$ contain the $K$ largest irrelevant teacher probabilities.
The unique threshold $t$ solves
\begin{equation}
\sum_{i\in R}(q_i-t)_+=\sum_{j\in S}(t-q_j)_+.
\label{eq:poolbalance}
\end{equation}
Define $A=\{i\in R:q_i>t\}$ and $B=\{j\in S:q_j<t\}$.
An attaining distribution is
\begin{equation}
p_i^*=\begin{cases}
\min(q_i,t),&i\in R,\\
\max(q_i,t),&i\in S,\\
q_i,&i\notin R\cup S.
\end{cases}
\end{equation}
Its pooling level and divergence are
\begin{align}
t&=\frac{\sum_{i\in A}q_i+\sum_{j\in B}q_j}{|A|+|B|},\\
C_K&=\sum_{i\in A}q_i\log\frac{q_i}{t}
+\sum_{j\in B}q_j\log\frac{q_j}{t}.
\label{eq:exactcost}
\end{align}
Coordinates equal to $t$ can join either active set without changing the result.
Zero-probability selected candidates can receive positive mass and contribute zero directly to KL.
\end{theorem}

\begin{proof}
For a fixed $K$-element set $S$, failure requires $p_j\geq p_i$ for all $j\in S$ and $i\in R$.
The complete event is the union over such sets.

First consider which irrelevant candidates to select.
Suppose $u\in S$, $v\in N\setminus S$, and $q_v\geq q_u$.
Take any feasible finite-cost distribution $p$.
If $p_v\geq p_u$, replacing $u$ by $v$ preserves feasibility without changing $p$.
Otherwise, swap $p_u$ and $p_v$.
The swapped distribution is feasible for the replacement set.
Its cost change is
\begin{equation}
(q_v-q_u)\log(p_v/p_u)\leq0.
\end{equation}
The boundary convention extends this argument to zero coordinates.
Repeated exchanges select the $K$ largest irrelevant teacher probabilities without increasing the minimum cost.

For this fixed set, introduce $t$ with constraints $p_i\leq t$ for relevant coordinates.
Selected irrelevant coordinates satisfy $p_j\geq t$.
Up to a constant, the objective is the convex function $-\sum_iq_i\log p_i$.
The proposed projection preserves total mass by Equation~\ref{eq:poolbalance}.
It also satisfies every order constraint.

Use simplex multiplier one.
For an active relevant coordinate, use multiplier $\mu_i=q_i/t-1\geq0$ on $p_i-t\leq0$.
For an active selected irrelevant coordinate, use $\nu_j=1-q_j/t\geq0$ on $t-p_j\leq0$.
All inactive order multipliers vanish.
Stationarity in each positive coordinate follows directly.
Stationarity in $t$ requires
\begin{equation}
\sum_{i\in A}\mu_i=\sum_{j\in B}\nu_j.
\end{equation}
This equality is the pooled-mean identity.
For an unpooled zero coordinate, nonnegativity multiplier one satisfies boundary stationarity.
Complementarity holds.
These convex optimality conditions prove minimality.

The balance function is continuous.
It is positive at the smallest selected irrelevant value and negative at the largest relevant value.
Between those bounds, at least one active coordinate makes it strictly decreasing.
Thus its root exists and is unique.
Substituting the projection into KL gives Equation~\ref{eq:exactcost}.
\end{proof}

\begin{corollary}
If $D(q\|p)<C_K(q,R)$, the student retains a relevant item within the first $K$ positions.
The open radius is optimal.
\end{corollary}
\begin{proof}
The inequality excludes the failure set by the definition of $C_K$.
The attaining projection reaches failure at equality under pessimistic ties.
Strict failures can approach that projection arbitrarily closely.
Thus no larger open KL radius works in general.
\end{proof}

\subsection{R@1 and boundary behavior}
For $K=1$, let $b$ be the largest irrelevant teacher probability.
Start a pool with $b$.
Add relevant probabilities in decreasing order while the next value exceeds the current mean.
If $m$ relevant coordinates join, then
\begin{align}
t&=\frac{b+\sum_{i=1}^m a_i}{m+1},\\
C_1&=\sum_{i=1}^m a_i\log(a_i/t)+b\log(b/t).
\end{align}
Adding a relevant value raises the mean without exceeding that value.
Previously active coordinates therefore remain active.
This proves the one-pass pooling procedure.

The R@1 cost is $O(n+r\log r)$.
General $K$ first selects the top $K$ irrelevant values.
The balance function is linear between its probability breakpoints.
The implementation scans active sets with cost $O((r+K)^2)$ after the gallery scan.
Relevant sets contain five captions or one image in the source retrieval tasks.

For $q=(.35,.30,.25,.10)$ and $R=\{1,2,3\}$, the exact radius is approximately $.08083268$ nats.
The two-coordinate radius is approximately $.07354844$ nats.
This deterministic example illustrates the contribution of multiple relevant candidates.
It is not an empirical retrieval result.

Whenever failure is possible, $C_K\leq\log(K+1)$.
To see this, select any $K$ irrelevant indices and define
\begin{equation}
p=\frac{q+\sum_{j\in S}e_j}{K+1}.
\end{equation}
Each selected irrelevant probability is at least every relevant probability.
Also $p_i\geq q_i/(K+1)$, which bounds the divergence.
A point mass on one relevant candidate attains the bound.

\subsection{An exact score radius}
Let $z$ and $z+d$ be teacher and student scores on the same gallery.
Let $z_{(K,N)}$ be the $K$th largest irrelevant score.
Define
\begin{align}
g_K&=\max_{i\in R}z_i-z_{(K,N)},\\
\operatorname{osc}(d)&=\max_i d_i-\min_i d_i\\
&=2\min_c\|d-c\mathbf1\|_\infty.\nonumber
\end{align}

\begin{theorem}
For a strictly correct teacher, $\operatorname{osc}(d)<g_K$ preserves a relevant item within the first $K$ positions.
The exact failure radius under the centered infinity norm is $g_K/2$.
\end{theorem}
\begin{proof}
Choose a teacher-best relevant candidate $i^*$.
Every irrelevant $j$ with $z_j\leq z_{(K,N)}$ satisfies
\begin{equation}
(z_{i^*}+d_{i^*})-(z_j+d_j)
\geq g_K-\operatorname{osc}(d)>0.
\end{equation}
At most $K-1$ irrelevant candidates lie above that threshold.
This proves retention.

For sharpness, decrease every relevant score by $g_K/2$.
Increase the $K$ largest irrelevant scores by $g_K/2$.
Leave all other scores unchanged.
The selected irrelevant candidates now tie or exceed every relevant candidate.
The oscillation is exactly $g_K$.
\end{proof}

This radius removes harmless common score shifts.
It improves the uncentered sufficient condition $2\|d\|_\infty<g_K$.
Consistent temperature scaling preserves the comparison.

For $q=\operatorname{softmax}(z/T)$ and $p=\operatorname{softmax}((z+d)/T)$,
\begin{equation}
D(q\|p)=\log\mathbb E_q e^{d/T}-\mathbb E_q[d/T].
\end{equation}
Hoeffding's bounded-variable inequality gives
\begin{equation}
D(q\|p)\leq\frac{\operatorname{osc}(d)^2}{8T^2}.
\end{equation}
Neither sufficient certificate dominates the other on every query.
Their union remains a valid certificate.

\subsection{A total KL budget}
Consider teacher-correct queries with positive finite thresholds $C_i$.
Suppose their full-gallery divergences sum to at most $B$.
Sort thresholds increasingly as $C_{(1)}\leq\cdots\leq C_{(m)}$.
The number of lost queries cannot exceed
\begin{equation}
\max\left\{h:\sum_{i=1}^hC_{(i)}\leq B\right\}.
\end{equation}
Each lost query costs at least its failure threshold.
Buying the cheapest thresholds therefore gives the largest possible failure count.
The bound is sharp when query distributions can vary independently.
Theorem~\ref{thm:exactprojection} supplies an attaining failure distribution for every threshold.
A shared embedding model can impose additional constraints.
The bound remains valid there, but need not be attainable.

\subsection{Scope and numerical safeguards}
The certificate uses the actual query, candidate gallery, relevance set, and temperature.
It cannot replace any of these with a training surrogate.
Two counterexamples make the boundary explicit.

First, take teacher scores $(2,1,0)$ with the first candidate relevant.
Train only on candidates one and two.
Student scores $(2,1,10)$ give zero KL on that training subgallery.
The full gallery nevertheless retrieves the third, irrelevant candidate.

Second, give 100 queries teacher probabilities $(.51,.49)$, with the first candidate relevant.
Leave 99 unchanged and switch one student to $(.49,.51)$.
Average KL can fall below the individual threshold while one query fails.
An average divergence cannot replace each query's divergence in the individual certificate.

The implementation rejects unsafe overflow rather than returning a passing certificate.
It uses the conservative comparison
\begin{equation}
D+10^{-12}+10^{-10}|C_K|<C_K.
\end{equation}
These tolerances do not constitute interval arithmetic.
The reported quantities remain floating-point diagnostics.
Benchmark ranking can use a fixed index tie rule.
Certificate correctness uses pessimistic ties and must remain separately labeled.

Any empirical coverage report must state its denominator.
Useful denominators include all queries, strictly correct teacher queries, and actually retained teacher-correct queries.
Frozen states must remain visible when interpreting coverage.
The present proofs do not assert a measured coverage gain.
